% Exact local bifurcation result, obtained after the fractional radial scan.
% The parent may adapt labels to its local-radius section.
\section{A proved failure of fractional global monotonicity}
\label{rwloc:turn:sec:main}

The local sign threshold does not extend directly to a classification of
radial monotonicity on the whole unit disk. There are fractional orders
for which $\eta_{a,b}(\rho)=\cos\theta_{a,b}(\rho)/\rho$ first
increases and then decreases. This is an analytic theorem, supported by
a finite exact certificate for its hypotheses.

\subsection{The next local coefficient}

Put $p_n=H_{n-1}^{(b)}/n^a$ and
\[
 \mu=\frac{p_3}{2p_2},\qquad
 K=-\frac{4p_3\mu^2-4p_4\mu+p_5}{2p_2}.
\]
Thus $\mu=\eta_{a,b}(0)$ and $K$ is the previously identified
quadratic coefficient. The next coefficient is
\begin{equation}\label{rwloc:turn:eq:quartic}
 Q=-\frac{(8p_3\mu-4p_4)K
              +8p_4\mu^3-12p_5\mu^2+6p_6\mu-p_7}{2p_2},
\end{equation}
and the convergent local expansion is
\begin{equation}\label{rwloc:turn:eq:etaexp}
 \eta_{a,b}(\rho)=\mu+K\rho^2+Q\rho^4+O(\rho^6).
\end{equation}
The coefficients and remainder are real analytic in $a,b$ locally
within the parameter range under consideration.

To check \eqref{rwloc:turn:eq:quartic}, write $c=\rho\eta$ and use
$\sin(n\theta)=\sin\theta\,U_{n-1}(c)$. Division of the
imaginary-part equation by $\rho^3\sin\theta$ gives
\begin{align*}
 0={}&2p_2\eta-p_3
   +\rho^2(4p_3\eta^2-4p_4\eta+p_5)\\
   &+\rho^4(8p_4\eta^3-12p_5\eta^2+6p_6\eta-p_7)
   +O(\rho^6).
\end{align*}
Every omitted mode has order at least $\rho^6$. This equation is
analytic in $a,b,t=\rho^2,\eta$ near $t=0$; it follows either
from the convergent disk series or from the parity of each Chebyshev
term. Its derivative in $\eta$ at zero is $2p_2>0$. The analytic
implicit-function theorem therefore gives \eqref{rwloc:turn:eq:etaexp},
and substitution proves the formula for $Q$.

\begin{theorem}[Emergence of a strict radial maximum]
\label{rwloc:turn:thm:maximum}
Fix $b=1/2$. There is a unique zero $a_\star$ of $K(a,1/2)$ in
the rational interval
\begin{equation}\label{rwloc:turn:eq:abracket}
 1.31298577512189<a_\star<1.31298577512190.
\end{equation}
There exist $\delta>0$ and $\rho_0>0$ such that, for every
$a\in(a_\star-\delta,a_\star)$, the function
$\eta_{a,1/2}(\rho)$ has exactly one critical point in
$(0,\rho_0)$. It is a strict local maximum: the function increases
before it and decreases after it within that interval.
Writing this point as $\rho_{\max}(a)$, one has
\begin{equation}\label{rwloc:turn:eq:maximumscale}
 \rho_{\max}(a)
 =\left(\frac{\partial_aK(a_\star,1/2)}
                    {2Q(a_\star,1/2)}\right)^{1/2}
       (a_\star-a)^{1/2}
       +O((a_\star-a)^{3/2}).
\end{equation}
The positive prefactor is rigorously enclosed by
\[
 44.0912731287<
 \left(\frac{\partial_aK(a_\star,1/2)}
                    {2Q(a_\star,1/2)}\right)^{1/2}
 <44.0912738633.
\]
In particular these fractional normalized-radius curves are not
monotone on $(0,1]$, despite having $K>0$.
\end{theorem}

\begin{proof}
The accompanying standard-library verifier
\src{certify_fractional_turning.py} uses only integer and rational
arithmetic. Let $I$ be the closed rational interval in
\eqref{rwloc:turn:eq:abracket}. Its output proves
\[
 K(\inf I,1/2)>8.16\cdot10^{-17},\qquad
 K(\sup I,1/2)<-8.72\cdot10^{-17},
\]
and, throughout $I$,
\begin{align*}
 -0.016895831331&<\partial_aK(a,1/2)<-0.016895831329,\\
 -4.345546\cdot10^{-6}&<Q(a,1/2)<-4.345545\cdot10^{-6}.
\end{align*}
These signs alone give existence and uniqueness of $a_\star$ in
$I$, independently of the more general local-threshold theorem.
That theorem also identifies it as the unique threshold in $(1,2)$.

Here is the complete arithmetic basis of the certificate. For
$2\leq n\leq7$, it encloses $\log n$ by the first $120$ terms of
\[
 \log n=2\sum_{j\geq0}\frac{z^{2j+1}}{2j+1},
 \qquad z=\frac{n-1}{n+1},
\]
with positive tail bounded by
$2z^{241}/(241(1-z^2))$. For a nonnegative rational $x$, the first
$60$ terms after the constant term of $e^x$ have tail at most
\[
 \frac{x^{61}}{61!}\frac1{1-x/62};
\]
all arguments used here are less than three. Monotonicity encloses
$e^{a\log n}$ at both ends of the interval arguments, and reciprocation
encloses $n^{-a}$. The factors $j^{-1/2}$ in the harmonic numbers are
bounded by exact integer-square-root operations. All arithmetic
operations round outwards to multiples of $2^{-128}$. Substitution
in the formulas for $K,Q$ and their elementary derivatives proves the
displayed signs. Exact rational endpoints, including the tighter
prefactor enclosure, are recorded in
\src{fractional_turning_certificate.json}. The written tail bounds
and integer inequalities make this a finite exact proof certificate;
no floating-point sign decision is used.

Write $\eta_{a,1/2}(\rho)=\Phi(a,t)$, $t=\rho^2$, using the
analytic function constructed above. Set $U(a,t)=\partial_t\Phi(a,t)$.
Then
\[
 U(a_\star,0)=0,\qquad
 \partial_tU(a_\star,0)=2Q(a_\star,1/2)<0,\qquad
 \partial_aU(a_\star,0)=\partial_aK(a_\star,1/2)<0.
\]
The analytic implicit-function theorem gives a unique local solution
$t=t(a)$ of $U(a,t)=0$, with
\[
 t(a)=\frac{\partial_aK(a_\star,1/2)}{2Q(a_\star,1/2)}
          (a_\star-a)+O((a_\star-a)^2).
\]
It is positive for $a<a_\star$ sufficiently close. On a sufficiently
small fixed rectangle, $\partial_tU<0$; at its upper $t$ boundary,
$U<0$, while $U(a,0)=K(a,1/2)>0$ for $a<a_\star$. Thus there
is exactly one such positive critical point and the derivative
$\partial_\rho\eta=2\rho U(a,\rho^2)$ changes from positive
to negative there. At the critical point its second derivative is
$4t\partial_tU<0$. Taking the positive square root of the expansion
for $t(a)$ proves \eqref{rwloc:turn:eq:maximumscale}.
\end{proof}

\subsection{Exploration beyond the existence theorem}

A separate 70-digit diagnostic evaluates the defining disk series and its
derivative at $b=1/4,1/2,3/4$, at the local threshold and at offsets
$-0.02,-10^{-4},+0.02$, for radii $0.3,0.7,0.9,0.98$. For each
$b$, the offset $-10^{-4}$ gives a positive derivative at $0.3$ and
a negative derivative at $0.98$. Thus the same phenomenon appears
numerically at all three sampled inner orders; the theorem above proves
it at $b=1/2$.

For the exact rational pair $a=21/16$, $b=1/2$, the diagnostic finds
\[
 \eta'(0.9)=1.2937181281564777\cdot10^{-6},\qquad
 \eta'(0.98)=-2.6349934635137804\cdot10^{-6}.
\]
These two particular signs are numerical observations, not an extra
certified instance of the theorem's unspecified interval in $a$.
The disk-series truncation errors for $F$ and $zF'$ are bounded
analytically by $10^{-50}$, but rounding is not interval-enclosed.
The existence theorem relies instead on the exact finite coefficient
certificate above. Extending this argument to other $b$ requires
determining the sign of $Q(a_\star(b),b)$. The complementary
minimum theorem below proves that this sign is not always negative.

% Continuation of fractional_turning.tex; uses its definitions of K and Q.
\subsection{A complementary minimum and a change of quartic sign}
\label{rwloc:turn:sec:minimum}

The sign of the quartic coefficient along the local threshold also
changes. In particular, the negative sign at $b=1/2$ cannot extend to
every $0<b<1$. The following exact certificate exhibits the opposite
local bifurcation at a rational inner order close to one.

\begin{theorem}[Emergence of a strict radial minimum]
\label{rwloc:turn:thm:minimum}
Set $b_0=999/1000$. The unique local threshold $a_0=a_\star(b_0)$
satisfies the rational enclosure
\begin{equation}\label{rwloc:turn:eq:minbracket}
 1.00057049936559087278<a_0<1.00057049936559087279.
\end{equation}
Moreover,
\[
 Q(a_0,b_0)>0,\qquad \partial_a K(a_0,b_0)<0.
\]
There exist $\delta>0$ and $\rho_0>0$ such that, for every
$a\in(a_0,a_0+\delta)$, the function $\eta_{a,b_0}(\rho)$
has exactly one critical point in $(0,\rho_0)$. It is a strict local
minimum: the function decreases before it and increases after it
within this interval. Its location satisfies
\begin{equation}\label{rwloc:turn:eq:minimumscale}
 \rho_{\min}(a)
 =\left(\frac{-\partial_a K(a_0,b_0)}{2Q(a_0,b_0)}\right)^{1/2}
       (a-a_0)^{1/2}+O((a-a_0)^{3/2}),
\end{equation}
where the prefactor has the certified enclosure
\[
 17076.2396129683<
 \left(\frac{-\partial_a K(a_0,b_0)}{2Q(a_0,b_0)}\right)^{1/2}
 <17076.2396548521.
\]
Thus fractional normalized-radius curves can also fail global
monotonicity while their initial quadratic coefficient is negative.
The theorem asserts a sufficiently small neighborhood; it gives no
numerical lower bound on $\delta$.
\end{theorem}

\begin{proof}
Let $I$ be the closed rational interval in
\eqref{rwloc:turn:eq:minbracket}. The exact interval verifier
\src{certify_fractional_minimum.py} proves
\[
 K(\inf I,b_0)>4.08027\cdot10^{-23},\qquad
 K(\sup I,b_0)<-1.29944\cdot10^{-22}.
\]
Throughout $I$, it also proves
\begin{align*}
 2.92779195\cdot10^{-11}
     &<Q(a,b_0)<2.92779198\cdot10^{-11},\\
 -0.017074763285092
     &<\partial_a K(a,b_0)<-0.017074763285091.
\end{align*}
The arithmetic uses the same outward dyadic intervals and rational
logarithm/exponential remainder bounds as the maximum certificate.
For this rational inner order, each $j^{-b_0}$ is enclosed as the
reciprocal of $e^{b_0\log j}$, for $1\leq j\leq6$; each
$n^{-a}$ is enclosed in the same way, for $2\leq n\leq7$.
All exponential arguments are less than three. Substitution into
the finite formulas for $K,Q$ and $\partial_aK$ establishes the
signs. Exact rational endpoints and the sharper prefactor interval
are in \src{fractional_minimum_certificate.json}. Thus both the
small positive quartic coefficient and the much smaller endpoint
values of $K$ are certified without floating-point arithmetic.

Write $\eta_{a,b_0}(\rho)=\Phi(a,t)$ with $t=\rho^2$, and
let $U=\partial_t\Phi$. At $(a_0,0)$ one has
\[
 U=0,\qquad U_t=2Q(a_0,b_0)>0,\qquad
 U_a=\partial_aK(a_0,b_0)<0.
\]
The analytic implicit-function theorem gives a unique local zero
$t=t(a)$ with
\[
 t(a)=\frac{-\partial_aK(a_0,b_0)}{2Q(a_0,b_0)}
                (a-a_0)+O((a-a_0)^2).
\]
It is positive for $a>a_0$ sufficiently close. On a sufficiently
small fixed rectangle, $U_t>0$; hence $U$ has exactly one zero in
the indicated positive $t$ interval and changes from negative to
positive there. Since $\eta_\rho=2\rho U$, this is the asserted
minimum, with $\eta_{\rho\rho}=4tU_t>0$ at the critical point.
Taking its positive square root proves
\eqref{rwloc:turn:eq:minimumscale}.
\end{proof}

\begin{corollary}[A zero of the quartic coefficient along the threshold]
\label{rwloc:turn:cor:quarticzero}
There exists $b_\dagger\in(1/2,999/1000)$ for which
\[
 Q(a_\star(b_\dagger),b_\dagger)=0.
\]
\end{corollary}

\begin{proof}
The threshold equation has nonzero $a$ derivative at every
$a_\star(b)$, so $a_\star(b)$ is real analytic on $(0,1)$.
The finite formula for $Q$ then makes
$q(b)=Q(a_\star(b),b)$ real analytic there. The exact maximum and
minimum certificates give $q(1/2)<0<q(999/1000)$, so the
intermediate value theorem applies.
\end{proof}

The finite-coefficient diagnostic
\src{locate_quartic_sign_transition.py}, run at 100 decimal
digits, locates one apparent sign change at
\[
 b_\dagger\approx0.9974937898734204210746055931,\qquad
 a_\star(b_\dagger)\approx1.0014301809614951381293517294.
\]
These displayed decimal locations are diagnostic, not certified
enclosures. No uniqueness of $b_\dagger$ is claimed. The 11-point
finite-coefficient sweep has negative $q(b)$ at every requested
sample from $0.001$ through $0.99$, and positive $q(0.999)$;
the latter sign is proved by the independent exact certificate.
Determining the number of zeros of $q$ and the higher-order radial
behavior at its zeros remains a concrete research problem.
